\uuid{QKFt}
\exo7id{2732}
\titre{exo7 2732}
\auteur{tumpach}
\organisation{exo7}
\datecreate{2009-10-25}
\isIndication{false}
\isCorrection{false}
\chapitre{Déterminant, système linéaire}
\sousChapitre{Système linéaire, rang}
\cpge{mpsi-mp2i.calcul-algebrique-trigonometrie.b, pcsi.calcul-algebrique-trigonometrie.b, ptsi.calcul-algebrique-trigonometrie.b, mpsi-mp2i.calcul-matriciel-systemes.c, pcsi.calcul-matriciel-systemes-lineaires.c, ptsi.calcul-matriciel-systemes-lineaires.c, mpsi-mp2i.matrices-applications-lineaires.c, pcsi.matrices-applications-lineaires.d, ptsi.matrices-applications-lineaires.d, bcpst1.systemes-lineaires.a, bcpst1.systemes-lineaires.b, bcpst1.systemes-lineaires.c}
\module{Algèbre}
\niveau{L2}
\difficulte{}

\contenu{
\question{
En appliquant l'algorithme de Gauss, r\'esoudre le syst\`eme lin\'eaire suivant~:
$$
\mathcal{S} ~: \left\{
\begin{array}{ccccccccc} 
2x_1 & + & 4x_2 & - & 6x_3 & - & 2x_4 & = & 2\\
3x_1 & + & 6x_2 & - & 7x_3 & + & 4x_4 &= & 2\\
5x_1 & + & 10 x_2 & - & 11 x_3 & + & 6x_4 &   = & 3
\end{array}
 \right.. 
$$
}
}
